\documentclass[12pt]{article}
\usepackage[T1]{fontenc}
\usepackage[T1]{polski}
\usepackage[utf8]{inputenc}
\newcommand{\BibTeX}{{\sc Bib}\TeX} 
\usepackage{graphicx}
\usepackage{amsfonts}
\usepackage{float}
\usepackage{amsmath}
\usepackage{subcaption}

\setlength{\textheight}{21cm}

\title{{\bf Zadanie nr 3 - Splot, filtracja i korelacja sygnałów}\linebreak
Cyfrowe Przetwarzanie Sygnałów}
\author{Maciej Miazek 251589 \and Remigiusz Tomecki 251652}
\date{14.05.2026 r.} 

\begin{document}
\clearpage\maketitle
\thispagestyle{empty}
\newpage
\setcounter{page}{1}

\section{Cel zadania}
Celem niniejszego zadania było zapoznanie się z fundamentalnymi operacjami cyfrowego przetwarzania sygnałów: splotem dyskretnym oraz korelacją wzajemną. Ważnym elementem ćwiczenia było praktyczne wykorzystanie twierdzenia o splocie do projektowania filtrów o skończonej odpowiedzi impulsowej (SOI) metodą okien oraz przeprowadzenie symulacji korelacyjnego czujnika odległości.

\section{Wstęp teoretyczny}
Podstawową operacją wykorzystywaną w filtracji cyfrowej jest splot dyskretny dwóch sygnałów $h$ oraz $x$, zdefiniowany wzorem:
\begin{equation}
(h*x)(n) = \sum_{k=0}^{M-1} h(k)x(n-k)
\end{equation}
W procesie projektowania filtrów metodą okien, idealna (nieskończona) odpowiedź impulsowa filtru dolnoprzepustowego jest mnożona przez funkcję okna $w(n)$, co pozwala na uzyskanie skończonej liczby współczynników $M$. Zastosowanie okien innych niż prostokątne (np. Hamminga, Hanninga, Blackmana) pozwala na redukcję tętnień w paśmie przepustowym i zwiększenie tłumienia w paśmie zaporowym.

Korelacja wzajemna sygnałów $R_{hx}$ służy do wyznaczania podobieństwa między sygnałami i jest kluczowa w systemach radarowych do wyznaczania opóźnienia czasowego sygnału zwrotnego, co pozwala na obliczenie odległości $d$:
\begin{equation}
d = \frac{V \cdot t}{2}
\end{equation}
gdzie $V$ to prędkość rozchodzenia się fali, a $t$ to czas wyznaczony na podstawie przesunięcia maksimum funkcji korelacji.

\section{Opis implementacji}
Aplikacja została napisana w języku Python. W ramach realizacji zadania posługiwaliśmy się bibliotekami: \verb|Matplotlib| - do generowania wykresów, \verb|Numpy| - do operacji na wektorach oraz \verb|Tkinter| - do utworzenia interfejsu użytkownika.

Główna logika filtracji została zawarta w klasie \verb|FiltrSygnalu|:
\begin{itemize}
    \item \textbf{Metoda \texttt{splot(h, x)}}: Realizuje bezpośrednią implementację wzoru na splot dyskretny przy pomocy zagnieżdżonych pętli, co pozwala na przetwarzanie sygnałów o dowolnej długości.
    \item \textbf{Projektowanie filtrów}: Zaimplementowano funkcje generujące filtry dolno-, górno- i środkowoprzepustowe. Wykorzystano technikę modulacji (mnożenie przez $(-1)^n$ dla filtrów górnoprzepustowych oraz przez $2\sin(\pi n/2)$ dla środkowoprzepustowych).
    \item \textbf{Funkcje okien}: Zaimplementowano okna: prostokątne, Hamminga, Hanninga oraz Blackmana zgodnie ze wzorami (5-7) z instrukcji.
\end{itemize}

\section{Eksperymenty i wyniki}

\subsection{Eksperyment 1: Charakterystyki filtrów SOI}
\subsubsection{Założenia i przebieg}
Przeprowadzono analizę wpływu rzędu filtru $M$ oraz rodzaju zastosowanego okna na charakterystykę amplitudową filtru dolnoprzepustowego o częstotliwości odcięcia $f_0 = f_p / 8$.

\subsubsection{Rezultaty}
Zgodnie z przewidywaniami teoretycznymi, zwiększenie rzędu filtru $M$ spowodowało zwężenie pasma przejściowego. Zastosowanie okna Blackmana w porównaniu do prostokątnego znacząco zredukowało oscylacje w pasmie zaporowym (tętnienia), kosztem nieznacznego poszerzenia pasma przejściowego.

\subsection{Eksperyment 2: Korelacyjny czujnik odległości}
\subsubsection{Założenia i przebieg}
Zasymulowano działanie radaru wysyłającego sygnał okresowy. Sygnał zwrotny został opóźniony o zadaną liczbę próbek, co odpowiadało określonej odległości obiektu.

\subsubsection{Rezultaty}
Wyznaczenie maksimum funkcji korelacji wzajemnej pozwoliło na precyzyjne określenie opóźnienia sygnału. Dla sygnału opóźnionego o 50 próbek przy znanej częstotliwości próbkowania, algorytm poprawnie wyliczył dystans, co potwierdza skuteczność metody korelacyjnej w pomiarach odległości.

\section{Wnioski}
1. Operacja splotu jest kluczowa dla filtracji sygnałów w dziedzinie czasu. Jej złożoność obliczeniowa $O(NM)$ sugeruje stosowanie zoptymalizowanych bibliotek (jak \texttt{numpy.convolve}) w rozwiązaniach komercyjnych, jednak implementacja bezpośrednia pozwala najlepiej zrozumieć mechanizm działania.
2. Wybór funkcji okna stanowi kompromis między szerokością pasma przejściowego a tłumieniem w pasmie zaporowym. Okno Blackmana oferuje najlepsze tłumienie, co jest pożądane w precyzyjnej filtracji.
3. Analiza korelacyjna jest potężnym narzędziem do detekcji sygnałów w szumie i wyznaczania przesunięć czasowych, co znajduje bezpośrednie zastosowanie w systemach radiolokacyjnych.

\end{document}